\documentclass[10pt,a4paper]{amsart}
\usepackage[margin=19mm]{geometry}
\usepackage{amsmath,amssymb,mathtools}
\usepackage[scr=rsfs,cal=euler]{mathalfa}
\usepackage{xfrac}
\usepackage[unicode,hidelinks,bookmarks=false]{hyperref}
\newcommand{\ie}{i.e.\ }
\newcommand{\cf}{cf.\ }
\newcommand{\resp}{respectively\ }
\newcommand{\Subsection}{\subsection}
\providecommand{\nobreakdash}{\nobreak\hbox{-}}
\setlength{\emergencystretch}{2em}
\multlinegap0pt
\usepackage{bm}
\usepackage{bbm}
\usepackage{dsfont}
\usepackage{xspace}
\usepackage{cancel}
\usepackage{mathtools}
\usepackage{amsthm}
\makeatletter
\@ifundefined{theorem}{\newtheorem{theorem}{Theorem}}{}
\@ifundefined{lemma}{\newtheorem{lemma}[theorem]{Lemma}}{}
\@ifundefined{proposition}{\newtheorem{proposition}[theorem]{Proposition}}{}
\makeatother
\title[Regularity of solution maps of the generalized surface quasi]{Regularity of solution maps of the generalized surface quasi-geostrophic equations}
\author{Gerard Misio{\l}ek, Xuan-Truong Vu, Tsuyoshi Yoneda}
\date{Source: arXiv:2603.12944v2 (CC BY 4.0)}
\begin{document}
\maketitle

\noindent\emph{Source and licence.} Regularity of solution maps of the generalized surface quasi-geostrophic equations. Version arXiv:2603.12944v2 (CC BY 4.0), distributed under the Creative Commons Attribution 4.0 International licence, as recorded in the arXiv licence metadata for this exact version and shown on the abstract page https://arxiv.org/abs/2603.12944v2. Full rights notice: licenses/arxiv-2603.12944-CC-BY-4.0.txt.\par\smallskip

\noindent\textit{Editorial scope.} Counted unit: the Lemma whose source label is \texttt{None} in the article (source0.tex lines 661--667, source label \texttt{None}), delivered together with the article's own complete original proof of it (source0.tex lines 668--718, one \texttt{proof} environment of 51 lines). No mathematics is written, repaired or re-derived: statement and proof are the article's own text line for line. Required context retained uncounted: prop:AIFT\_Shnirelman (lines 1382--1388). Every definition, display and label of those lines is retained unchanged. Omitted from this package and not redistributed: 1--660, 719--1381, 1389--1531. Nothing inside a retained excerpt is omitted or altered: every retained body is delivered exactly as the article's lines are written, so no omission receipt of any kind accompanies this package. Citations: the retained excerpts carry no citation command at all, so no bibliography entry of the article is transcribed and no reference is redistributed. Reference adaptation: none was needed and none was made; every reference inside the delivered unit resolves inside it, so no unresolved reference or citation is produced. Printed slips kept literal: no printed word, symbol, hypothesis or number of the counted statement or of its proof was altered. Layout and class: the article's own class is \texttt{amsart} with packages including inputenc, geometry, graphicx, amsmath, amsfonts, amssymb, latexsym, amsthm, epstopdf, enumerate, appendix, xcolor, hyperref, chngcntr; the wrapper is the lane's pinned \texttt{amsart} editorial wrapper at 10pt on A4, which restates the theorem-like environments the retained text needs and adds the article's own macro definitions that the retained text needs (none), with extra wrapper packages (bm, bbm, dsfont, xspace, cancel, mathtools). The delivered numbering is the unit's own. Assets: no retained span contains a figure environment, a graphics-inclusion command, a PDF-page inclusion, a table or a TikZ picture; no image file is copied into this package, the package declares no asset dependency and no image file is redistributed with it. Licence: version arXiv:2603.12944v2 of this article is distributed under the Creative Commons Attribution 4.0 International licence, as recorded in the arXiv licence metadata for this exact version and shown on the abstract page https://arxiv.org/abs/2603.12944v2. A full rights notice is delivered at licenses/arxiv-2603.12944-CC-BY-4.0.txt. Characters: every retained excerpt is pure ASCII, so the delivered engine, XeLaTeX with its default Latin Modern text font, reports no missing character.
\begin{proposition}\label{prop:AIFT_Shnirelman}
  Let $E, F, G$ be complex Banach spaces, and let $\Phi: E \times F \to G$ be an analytic mapping defined in a neighborhood of a point $(x_0, y_0) \in E \times F$, with $\Phi(x_0, y_0) = z_0$. Suppose that the partial derivative $\frac{\partial \Phi}{\partial x} (x_0, y_0)$ is an invertible linear operator whose image is all of $G $. Then:  

1. There exists a radius $r > 0$ such that for any $y \in F$ and $z \in G$ satisfying $\|y - y_0\|_F \leq r$ and $\|z - z_0\|_G \leq r$, the equation $\Phi(x, y) = z$ has a unique solution $x(y, z)$ in a neighborhood of $x_0$.  

2. The function $x(y, z)$ is analytic with respect to $(y, z)$.  
\end{proposition}

 \begin{lemma}
  Let $\varepsilon>0$ be sufficiently small. If $v\in H_0^{s}$ with $\|v\|_s<\varepsilon$ then there is a unique (up to an additive constant) function $\varphi_{v}\in H^{s+1}$ such that the map 
     \begin{align}
     	x\rightarrow \xi_v(x)=x+v(x)+\nabla \varphi_{v}(x)
     \end{align}
 is an exact volume-preserving $H^s$ diffeomorphism. Furthermore, the map $v\mapsto v+\nabla  \varphi_{v}$ is analytic  as a map from $H_0^s$ to $H^s$.
 \end{lemma}

\begin{proof}
Since  the map $x\mapsto x+v(x)$ does not in general preserve the volume form $\mu$, we need to adjust it by adding a gradient term $\nabla \varphi_{v}(x)$ which can be done with the help of the Hodge decomposition.
    Call the adjusted map $\xi_v=e+v+\nabla \varphi_{v}$ satisfying $\det (J (\xi_v))=1$. So $\xi_v$ is an element of the subgroup $\mathcal{D}_{\mu, ex}^s$ near the identity and we have
     \begin{align}\label{eqn:det_der_gamma_v}
         \det\begin{pmatrix}
             1+\dfrac{\partial v_1}{\partial x_1}+ \dfrac{\partial^2 \varphi_{v}}{\partial x_1^2} &  \dfrac{\partial v_1}{\partial x_2}+ \dfrac{\partial^2 \varphi_{v}}{\partial x_1 \partial x_2}\\
             \dfrac{\partial v_2}{\partial x_1}+ \dfrac{\partial^2 \varphi_{v}}{\partial x_2 \partial x_1} &  1+\dfrac{\partial v_2}{\partial x_2}+ \dfrac{\partial^2 \varphi_{v}}{\partial x_2^2}
         \end{pmatrix}=1.
     \end{align}

      Expanding the \eqref{eqn:det_der_gamma_v} and rewriting it we obtain an equation for the Laplacian of $\varphi_v$, that is
      \begin{align}\label{eqn:contraction_map}
      	\Delta\varphi_{v}= P(D v,D^2\varphi_{v}).
      \end{align}
  where $P$ is a polynomial of degree $2$ in $Dv$ and $D^2\varphi$, namely 
 \begin{align}
    P\big(D v,D^2\varphi\big)= -\bigg(\frac{\partial v_1}{\partial x_1}+ \frac{\partial v_2}{\partial x_2}
     &+ \frac{\partial v_1}{\partial x_1} \frac{\partial v_2}{\partial x_2}+ \frac{\partial v_1}{\partial x_1} \frac{\partial^2 \varphi}{\partial x_2^2}+ \frac{\partial v_2}{\partial x_2} \frac{\partial^2 \varphi}{\partial x_1^2} + \frac{\partial^2\varphi}{\partial x_1^2}\frac{\partial^2\varphi}{\partial x_2^2} \bigg)\notag\\
     &+  \frac{\partial v_2}{\partial x_1} \frac{\partial v_1}{\partial x_2}  + \frac{\partial v_2}{\partial x_1}  \frac{\partial^2 \varphi}{\partial x_1 \partial x_2}+ \frac{\partial v_1}{\partial x_2} \frac{\partial^2 \varphi}{\partial x_2 \partial x_1} + \bigg(\frac{\partial^2 \varphi}{\partial x_1 \partial x_2}\bigg)^2.
 \end{align}
 Let $Q(v,\varphi_{v})=  	\Delta^{-1} P(D v,D^2\varphi_{v})$.
Observe that $Q$ is analytic in both variables as a map from $H_0^s\times H^{s+1}$ to $H^{s+1}$ and, moreover, $Q$ and its G\^ateaux derivatives satisfy the inequalities
\begin{align}\label{eqn:bound_Q}
    &\|Q(v,\varphi)\|_{s+1}\le C\big(\|v\|_s^2+\|\varphi\|_{s+1}^2\big) 
\end{align}
and
\begin{align}\label{eqn:bound_nabla_Q}
    &\|\partial_vQ(v,\varphi)\|_{s+1}+\|\partial_{\varphi} Q(v,\varphi)\|_{s+1}\le C(\|v\|_s+\|\varphi\|_{s+1}),
\end{align}
with some constant $C>0$. 

Let $\varepsilon<(2C)^{-1}$. For any $v$ in $H_0^s$ with $\|v\|_{s}<\varepsilon$, if $\|\varphi\|_{s+1}\le \varepsilon$  then from \eqref{eqn:bound_Q} we have 
 \begin{align}
    	\|Q(v,\varphi)\|_{s+1}\le C(\varepsilon^2+\varepsilon^2)=2C\varepsilon^2< \varepsilon
    \end{align}
which implies that $\varphi \rightarrow Q(v,\varphi)$ maps the $\varepsilon$-ball $B_{\varepsilon}^{s+1}$ centered at the origin in $H^{s+1}$ to itself. 
   
Furthermore, from \eqref{eqn:bound_nabla_Q}, for any $\varphi_1,\varphi_2$ in $B_{\varepsilon}^{s+1}$ we have 
%
    \begin{align}
    \|Q(v,\varphi_1)-Q(v,\varphi_2)\|_{s+1}&\le \max_{0\le \lambda\le 1}\|\nabla Q\big(v,\lambda\varphi_1+(1-\lambda)\varphi_2\big)\| \|\varphi_1-\varphi_2\|_{s+1}\notag\\
  %  &\le C(\varepsilon+\varepsilon) \|\varphi_1-\varphi_2\|_{s+1}\\
    &\le 2C\varepsilon \|\varphi_1-\varphi_2\| _{s+1}\\
    &<  \|\varphi_1-\varphi_2\| _{s+1}.\notag
\end{align}
%
    Thefore, there is a ball of sufficiently small radius such that for any $v$ belongs to this ball, there exists a unique solution $\varphi=\varphi_v$ of \eqref{eqn:contraction_map} by the contraction mapping argument. 
    
    Since $Q$ is analytic, by the Analytic Implicit Function Theorem \ref{prop:AIFT_Shnirelman}, there exists a unique analytic solution $\varphi_v$ which shows that the map $v\mapsto\xi_v-e$ is analytic  as a map from $H_0^s$ to $H^s$.

 \end{proof}

\end{document}
